\scanpage{111}
\mspage{106}
% archive sheet number: 110

\uline{Corollaire 2} \inserted{Sous les conditions \inserted{\unsure{équ.}} du Cor.\,1 (iii),} Si $\alpha,\beta\ne0$, $\alpha\vee\beta\ne1$, et posant $\unsure{\gamma|\sigma}=\unsure{C|\complement}(\alpha\vee\beta)\ne0$, $\alpha,\beta,\gamma$ forment une ``base'' de l'algèbre de Boole engendrée par $\alpha,\beta$. Sous les conditions du lemme 2, \inserted{et supposant $0<b<c<1$,} posant
\[
\alpha=b,\quad\beta=c\wedge Cb,\quad\gamma=Cc
\]
$\alpha,\beta,\gamma$ forment une base de l'algèbre de Boole engendrée par $b,c$.

\noindent\makebox[0pt][r]{\raisebox{-4\baselineskip}[0pt][0pt]{\rule{0.3pt}{5\baselineskip}\kern2pt\rule{0.3pt}{5\baselineskip}}\hspace{0.8em}}\uline{Question} Les conditions envisagées au \deleted{lemme 1} Cor.\,1 sur $\unsure{\Sigma|\bar\Sigma}$ sont elles tjrs vérifiées ? J'en doute fort ! Le sont elles si $\Sigma$ est fini${}^{*}$ ? J'en doute également.
% slanted
\margin{Non, \unsure{contrexemple|contre-exemple} avec $\operatorname{Card}\Sigma=10$}
Si on suppose ces conditions satisfaites, $\Sigma$ est-il nécessairement une algèbre de Boole ? L'est-il si on fait l'hypothèse similaire à celle du Cor.\,1 \inserted{(\unsure{iiii'|iiii|iv})}, pour \uline{$\unsure{\forall|\text{tt}}$} familles $(\alpha_i)_{\unsure{\nu|i}\in I}$ d'éléments mutuellement disjoints de $\unsure{I|\Sigma}$ (pas néc. finie).

\begin{quote}
\noindent\vrule\kern2pt\vrule\quad\begin{minipage}{0.86\linewidth}
Il faudrait absolument avoir le coeur net de ce qui en est, quand $\Sigma$ est défini via un ens. ordonné $L$, muni de la relation $\unsure{a|x}|\unsure{\circ|0|\bullet}|y$ sss $x<y$ ou $y<x$.
\end{minipage}
\end{quote}
